% ID: SEL-D0005-OB-E02
% ESTADO: cotejada. Segunda pasada visual completa y borrador compilado revisados.
% Fuente: Selectividad/sel_2019_matematicas_aplicadas/Reserva A/Exámenes Mat Aplicadas CCSS - A y B.pdf
% SHA-256: d28871f7f203cf1be9e479e9d642cdb6b1623704a8b76a2b57cfa2d12950f293
% Página: 2; opción B; ejercicio 2.

Unos productores de cereales realizan un estudio para conocer la posible
demanda de su producto. Concluyen que la función de demanda de dichos cereales
tiene la forma
\[
  D(x)=-200x^3+2\,100x^2-7\,200x+10\,000,
\]
para $0\leq x\leq 4$, donde $x$ es el precio en euros por kilogramo de
producto y $D(x)$ es la cantidad de kilogramos de cereales que los consumidores
están dispuestos a comprar a dicho precio $x$.
\begin{enumerate}[label=\alph*)]
  \item \textbf{(0.5 puntos)} \textquestiondown{}Cuál es la cantidad de
  cereales demandada si el precio es de 0.50 euros por kilogramo?
  \item \textbf{(2 puntos)} Calcule para qué precio se alcanza una demanda
  mínima del producto y determine dicha demanda.
\end{enumerate}
